\documentclass{beamer}

\title{Calibration-First Lesion Segmentation Update}
\author{}
\date{}

\begin{document}

\begin{frame}
  \titlepage
\end{frame}

\begin{frame}{Motivation}
  \begin{itemize}
    \item Shift the update from mask-only accuracy to reliable confidence.
    \item Treat calibration as a first-class requirement for lesion review.
    \item Support safer thresholding, triage, and downstream measurement.
  \end{itemize}
\end{frame}

\begin{frame}{Current Focus}
  \begin{itemize}
    \item Measure segmentation quality with Dice and boundary-sensitive checks.
    \item Measure probability quality with calibration curves and expected calibration error.
    \item Compare operating points before and after calibration-aware tuning.
  \end{itemize}
\end{frame}

\begin{frame}{Update Plan}
  \begin{itemize}
    \item Audit confidence distributions across lesion size, site, and image quality.
    \item Tune thresholds using validation-set calibration behavior.
    \item Report paired mask metrics and confidence reliability metrics.
  \end{itemize}
\end{frame}

\begin{frame}{Next Readout}
  \begin{itemize}
    \item Show qualitative examples where calibration changes decisions.
    \item Summarize tradeoffs between Dice, sensitivity, and false positives.
    \item Decide whether calibrated probabilities are ready for reviewer-facing outputs.
  \end{itemize}
\end{frame}

\end{document}
